\documentclass[conference]{IEEEtran}
\usepackage{cite}

\ifCLASSINFOpdf
\else
\fi

\usepackage{amsmath}
\usepackage{url}
\usepackage{graphicx}
\usepackage[hidelinks,bookmarks=false]{hyperref}
\usepackage{scalerel}
\usepackage{tikz}
\usepackage[utf8]{inputenc}
\usepackage[english]{babel}
\usepackage{siunitx}
\usepackage{float}
\usepackage[table]{xcolor}
\usepackage{booktabs}



\begin{document}

%%%%%%%%%%%%%%%%%%%%%%%%%%%%%%%%%%%%%%%%%%%%%%%%%%%%%%%%%%%%%%%%%%%%%%%%%%%%%
% Rename abstract and keywords
%%%%%%%%%%%%%%%%%%%%%%%%%%%%%%%%%%%%%%%%%%%%%%%%%%%%%%%%%%%%%%%%%%%%%%%%%%%%%
\renewcommand\IEEEkeywordsname{Keywords}

%%%%%%%%%%%%%%%%%%%%%%%%%%%%%%%%%%%%%%%%%%%%%%%%%%%%%%%%%%%%%%%%%%%%%%%%%%%%%


\title{DC Motor Angular Speed Control hps5130}

\author{
\IEEEauthorblockN{
Eric Castro Velásquez\IEEEauthorrefmark{1},
Esteban Segura Garcia\IEEEauthorrefmark{1},
Daniel Chavarria Garcia\IEEEauthorrefmark{1}
}
\vspace{2mm}
\IEEEauthorblockA{
\IEEEauthorrefmark{1}Escuela de Ingeniería Electrónica,
Instituto Tecnológico de Costa Rica (ITCR), 30101 Cartago, Costa Rica, \\
\{ecastro, essegura, dachavarria\}@estudiantec.cr
}
}

\maketitle

\begin{abstract}
This paper presents the design and validation.....
\end{abstract}

\begin{IEEEkeywords}
DC Motor, controller, PI, stabilization, overshoot.
\end{IEEEkeywords}

\IEEEpeerreviewmaketitle

\section{Introduction}

The control of angular velocity in motor of direct current, is a classic problem of automatic control with a big application in industrials systems and electric conditioning. A typical solution for this type of plant its the proportional-integrative regulator (PI), whose integrative action guarantee a zero steady-state error for inputs like step-type and the capacity of reject disturbances.\cite{ogata}

In this experiment, it is designed a PI regulator for the CD motor system \textit{hps5130}, modelling like a first order system, whose being identify via experimental. The design is performed for both continuous and discrete time with a sampling period of $T_s = \SI{5}{\milli\second}$, meeting the specifications of a settling time of less than 2\% for 300 ms and overshoot of less than 3\%. The results of both designs are verified by simulation in \textit{Simulink}, including the response to disturbances apply to the plant's input and output.\cite{guia}


\section{Theoretical Background}

\subsection{Motor CD as the control plant}

The motor's DC current can be modeled as a first-order linear system in the frequency domain, whose transfer function is related to the armature voltage applied to the output angular speed:
\begin{equation}
    G(s) = \frac{K}{s + p},
    \label{eq:modelo_motor}
\end{equation}

\noindent where $K$ is the static gain and $p$ is the dominant pole, both obtained through experimental identification. \cite{nise}


\subsection{PI regulator}

The PI regulator can be expressed in continuous time as:

\begin{equation}
    PI(s) = K_P + \frac{K_I}{s} = \frac{K_c(s - z_0)}{s},
    \label{eq:pi_continuo}
\end{equation}

\noindent where $K_P$ determines the system's response speed and $K_I$ eliminates the steady-state error for step inputs. In its transfer function form, the PI introduces a pole at the origin $s = 0$ and a zero at $s = z_0$, whose relative location with respect to the plant pole directly influences the closed-loop system dynamics.\cite{nise} 

\subsection{Root-locus method}

The root locus is a graphic technique that allows visualization of the displacement of the poles of the closed-loop in the complex plane $s$, in function of the controller gain. For a system with unity feedback, the poles of the closed-loop satisfy:

\begin{equation}
    1 + K_c \cdot G_{OL}(s) = 0,
    \label{eq:lazo_cerrado}
\end{equation}

\noindent where $G_{OL}(s)$ is the open loop transfer function. The design of the regulator consist in locate the poles of closed-loop in the region's of the plane $s$ that satisfy simultaneously the specifications of settling time and overshoot.\cite{nise} For a second-order system, these specifications impose restrictions on the real part and the angle of the dominant pole, expressed as:

\begin{equation}
    \sigma = \frac{4}{t_s},
    \label{eq:sigma}
\end{equation}

\begin{equation}
    \zeta = \frac{-\ln(M_p/100)}{\sqrt{\pi^2 + \ln^2(M_p/100)}},
    \label{eq:zeta}
\end{equation}

\noindent where $\sigma$ is the minimum real part required for the dominant pole, $t_s$ is the settling time to the 2\% and $\zeta$ is the damping coefficient corresponding to the maximum overshoot $M_p$ 

\subsection{Controller discretization}

To implement a PI regulator in a system, the controller in the continuous time has to be convert into discrete time. Using the backward Euler approximation, the complex variable $s$ is replaced by:  

\begin{equation}
    s \approx \frac{z - 1}{T_s \cdot z},
    \label{eq:euler_backward}
\end{equation}

\noindent which transforms the continuous PI into its discrete equivalent:

\begin{equation}
    PI(z) = \frac{b_1 z + b_0}{z + a_0},
    \label{eq:pi_discreto}
\end{equation}

\noindent where the coefficients $b_1$, $b_0$ y $a_0$ depends of $K_P$, $K_I$ and the sampling period $T_s$. The selection of $T_s$ must guarantee between 30 and 75 samples during the stabilization time for an adequate representation of the system dynamics.\cite{ogata}\cite{nise}

\subsection{Performance Criteria}

The design specifications are expressed in terms of three main metrics. 
The 2\% settling time ($t_s$) is defined as the time required for the 
response to remain within a $\pm 2\%$ band around the final value. The 
maximum percent overshoot ($M_p$) represents the largest percentage 
deviation of the response with respect to that final value. Finally, 
the steady-state error is the difference between the reference value 
and the final output value under a step input, which is zero when the 
controller contains an integrator in its structure \cite{ogata}.

\section{Methodology}

This section outlines the design and experimental validation process of 
three PI controllers for the angular velocity regulation of the 
\textit{hps5130} DC motor, in both continuous and discrete time.

\subsection{Plant Model Selection}

Multiple CSV datasets were collected from the \textit{hps5130} motor during 
the identification experiment. A statistical analysis based on the mean and 
standard deviation of the recorded signals was performed to determine the 
most representative dataset, which was subsequently used to obtain the 
empirical first-order plant model via the \textit{System Identification} 
toolbox in MATLAB, following the structure defined in \eqref{eq:modelo_motor}.

\subsection{PI Controller Design}

Three PI controllers were designed in both continuous and discrete time 
using the \textit{sisotool} root locus editor in MATLAB. Each design 
differed in the location of the real zero relative to the dominant plant 
pole, following the three placement criteria described in 
\eqref{eq:pi_continuo}. All designs were constrained to meet a 2\% settling 
time $t_s < \SI{300}{\milli\second}$ and a maximum percent overshoot 
$M_p < \SI{3}{\percent}$. The continuous-time controllers were verified 
in \textit{Simulink} under step input and disturbance rejection conditions, 
while the discrete-time equivalents were obtained at $T_s = \SI{5}{\milli\second}$ 
using the \texttt{c2d} function in MATLAB, but also, were verified in \textit{Simulink} under step input and disturbance rejection conditions.

\subsection{Experimental Validation}

Falta esta parte, se tendría que completar luego del miércoles

\section{Results}

This section presents the simulation and experimental results obtained 
for the three PI controllers designed in both continuous and discrete time.

\subsection{Continuous-Time Simulation Results}

\subsubsection{Step Response}


\begin{figure}[H]
    \centering
    \includegraphics[width=\columnwidth]{figures/continuous_step.png}
    \caption{Step response of the three continuous-time PI controllers 
    obtained in Simulink.}
    \label{fig:continuous_step}
\end{figure}



\subsubsection{Disturbance Rejection}


\begin{figure}[H]
    \centering
    \includegraphics[width=\columnwidth]{figures/continuous_disturbance.png}
    \caption{Disturbance rejection response of the three continuous-time 
    PI controllers obtained in Simulink.}
    \label{fig:continuous_disturbance}
\end{figure}



\subsection{Discrete-Time Simulation Results}

\subsubsection{Step Response}


\begin{figure}[H]
    \centering
    \includegraphics[width=\columnwidth]{figures/discrete_step.png}
    \caption{Step response of the best of three discrete-time PI controllers 
    obtained in Simulink.}
    \label{fig:discrete_step}
\end{figure}



\subsubsection{Disturbance Rejection}

\begin{figure}[H]
    \centering
    \includegraphics[width=\columnwidth]{figures/discrete_disturbance.png}
    \caption{Disturbance rejection response of the best of three discrete-time 
    PI controllers obtained in Simulink.}
    \label{fig:discrete_disturbance}
\end{figure}


\subsection{Continuous-Time Experimental Results}

\subsubsection{Step Response}

%
%\begin{figure}[H]
%    \centering
%    \includegraphics[width=\columnwidth]{figures/experimental_step.png}
%    \caption{Experimental step response of the implemented PI controllers 
%    on the \textit{kitprog059} platform on continuous-time.}
%   \label{fig:experimentalc_step}
%\end{figure}



\subsubsection{Disturbance Rejection}


%\begin{figure}[H]
%    \centering
%    \includegraphics[width=\columnwidth]{figures/experimental_disturbance.png}
%    \caption{Experimental disturbance rejection response of the implemented 
%    PI controllers on the \textit{kitprog059} platform on continuous-time.}
%    \label{fig:experimentalc_disturbance}
%\end{figure}


\subsection{Discrete-Time Experimental Results}

%\begin{figure}[H]
%    \centering
%    \includegraphics[width=\columnwidth]{figures/experimental_step.png}
%    \caption{Experimental step response of the implemented PI controllers 
%    on the \textit{kitprog059} platform on discrete-time.}
%    \label{fig:experimentald_step}
%\end{figure}



\subsubsection{Disturbance Rejection}


%\begin{figure}[H]
%    \centering
%    \includegraphics[width=\columnwidth]{figures/experimental_disturbance.png}
%    \caption{Experimental disturbance rejection response of the implemented 
%    PI controllers on the \textit{kitprog059} platform on discrete-time.}
%    \label{fig:experimentald_disturbance}
%\end{figure}

\subsection{Performance Summary}

\begin{table}[H]
    \centering
    \caption{Simulation performance metrics of the designed PI controllers.}
    \label{tab:performance}
    \begin{tabular}{ccccc}
        \hline
        \textbf{Controller} & \textbf{Domain} & 
        \textbf{$t_s$ (ms)} & \textbf{$M_p$ (\%)} & 
        \textbf{$e_{ss}$} \\
        \hline
        PI$_1$ & Continuous & 310 & 0 & 0 \\
        PI$_2$ & Continuous & 391 & 0 & 0 \\
        PI$_3$ & Continuous & & & \\
        \hline
        PI$_1$ & Discrete   & 265 & 0.49 & 0 \\
        PI$_2$ & Discrete   &370 & 0 & 0 \\
        PI$_3$ & Discrete   & & & \\
        \hline
    \end{tabular}
\end{table}

\begin{table}[H]
    \centering
    \caption{Experimental performance metrics of the designed PI controllers.}
    \label{tab:performance}
    \begin{tabular}{ccccc}
        \hline
        \textbf{Controller} & \textbf{Domain} & 
        \textbf{$t_s$ (ms)} & \textbf{$M_p$ (\%)} & 
        \textbf{$e_{ss}$} \\
        \hline
        PI$_1$ & Continuous & & & \\
        PI$_2$ & Continuous & & & \\
        PI$_3$ & Continuous & & & \\
        \hline
        PI$_1$ & Discrete   & & & \\
        PI$_2$ & Discrete   & & & \\
        PI$_3$ & Discrete   & & & \\
        \hline
    \end{tabular}
\end{table}


\section{Discussion}

\subsection{Statistical analysis}



Tables \ref{tab:velocity_results} and \ref{tab:current_results} summarize the parameter estimation results for both velocity and current models. In the case of velocity, the gain $K$ exhibits excellent repeatability, with a low coefficient of variation of 1.64\%, indicating a highly consistent estimation across repetitions. However, the time constant $\tau$ presents a significantly high variability (CV = 47.14\%), suggesting that this parameter is strongly affected by experimental conditions, measurement noise, or variations in the transient response.

For the current model, the behavior is notably different. The gain $K$ shows moderate variability (CV = 13.93\%), which can still be considered acceptable but indicates some dispersion between measurements. On the other hand, the time constant $\tau$ demonstrates excellent repeatability, with a very low coefficient of variation of 4.61\%, indicating stable and reliable estimation.

Overall, these results indicate that the estimation quality depends on both the measured variable and the parameter of interest. While the gain in velocity and the time constant in current are highly reliable, the time constant in velocity requires further analysis or additional repetitions to reduce uncertainty.

\begin{table}[H]
\centering
\caption{Experimental results for velocity model identification}
\label{tab:velocity_results}
\resizebox{\columnwidth}{!}{
\begin{tabular}{c c c c c}
\hline
Parameter & Rep. 1 & Rep. 2 & Rep. 3 & Statistics \\
\hline
$K$ & 1.3342 & 1.3708 & 1.3746 & $\mu=1.3598$, $s=0.0223$, $CV=1.64\%$ \\
$\tau$ (s) & 0.2301 & 0.1099 & 0.1071 & $\mu=0.1490$, $s=0.0703$, $CV=47.14\%$ \\
\hline
\end{tabular}
}
\end{table}


\begin{table}[H]
\centering
\caption{Experimental results for current model identification}
\label{tab:current_results}
\resizebox{\columnwidth}{!}{
\begin{tabular}{c c c c c}
\hline
Parameter & Rep. 1 & Rep. 2 & Rep. 3 & Statistics \\
\hline
$K$ & 0.1644 & 0.1326 & 0.1282 & $\mu=0.1418$, $s=0.0197$, $CV=13.93\%$ \\
$\tau$ (s) & 0.1090 & 0.1081 & 0.1174 & $\mu=0.1115$, $s=0.0051$, $CV=4.61\%$ \\
\hline
\end{tabular}
}
\end{table}



Table \ref{tab:rmse_results} presents the RMSE obtained during both training and validation for the three identified transfer functions. The model trained with Rep.~1 shows the lowest training RMSE (0.0952), indicating a very good fit to its own dataset. However, its validation RMSE is significantly higher (0.4089), suggesting poor generalization and possible overfitting.

In contrast, the models trained with Rep.~2 and Rep.~3 exhibit slightly higher training RMSE values (0.1122 and 0.1114, respectively), but substantially lower validation errors. Among them, the model trained with Rep.~3 achieves the lowest validation RMSE (0.2113), closely followed by Rep.~2 (0.2220). This indicates that both models generalize better to unseen data, with Rep.~3 providing the most accurate prediction overall.

According to the comparison criteria between training and validation RMSE, the models trained with Rep.~2 and Rep.~3 demonstrate adequate generalization, while the model trained with Rep.~1 fails to generalize properly.

\begin{table}[H]
\centering
\caption{Training and validation RMSE for the identified transfer functions}
\label{tab:rmse_results}
\resizebox{\columnwidth}{!}{
\begin{tabular}{c c c c}
\hline
Training & Validation & RMSE (fit) & RMSE (validation) \\
\hline
Rep. 1 & Rep. 2, 3 & 0.0952 & 0.4089 \\
Rep. 2 & Rep. 1, 3 & 0.1122 & 0.2220 \\
Rep. 3 & Rep. 1, 2 & 0.1114 & 0.2113 \\
\hline
\end{tabular}
}
\end{table}

A comprehensive evaluation considering statistical metrics and RMSE values allows selecting the most appropriate transfer function. From the statistical analysis, the gain $K$ in velocity exhibits excellent repeatability, whereas the time constant $\tau$ shows high variability , indicating uncertainty in its estimation. 

Between these two, the model trained with Rep.~3 yields the lowest validation RMSE, making it the most reliable representation of the system dynamics. Additionally, its performance is consistent with the acceptable variability observed in the gain and the more stable behavior seen in current-based parameters.

Therefore, the transfer function obtained from Rep.~3 is selected as the final model, as it provides the best compromise between fitting accuracy and generalization capability.
\subsection{Experimental versus simulation results}

\section{Conclusions}


%%%%%%%%%%%%%%%%%%%%%%%%%%%%%%%%%%%%%%%%%%%%%%%%%%%%%%%%%%%%%%%%%%%%%%%%%%%%%
\begin{thebibliography}{99}

\bibitem{ogata}
K. Ogata, \textit{Modern Control Engineering}, 5th ed. 
Upper Saddle River, NJ: Prentice Hall, 2010.

\bibitem{nise}
N. S. Nise, \textit{Control Systems Engineering}, 7th ed. 
Hoboken, NJ: Wiley, 2015.

\bibitem{guia}
E. Interiano and A. Ruiz, ``Proyecto Corto: Control de velocidad 
angular del motor CD hps5130,'' Laboratorio de Control Automático, 
EL-5409, Escuela de Ingeniería en Electrónica, Instituto Tecnológico 
de Costa Rica, San José, Costa Rica, Mar. 2025.

\end{thebibliography}
%%%%%%%%%%%%%%%%%%%%%%%%%%%%%%%%%%%%%%%%%%%%%%%%%%%%%%%%%%%%%%%%%%%%%%%%%%%%%
\end{document}


